\documentclass[10pt,a4paper]{amsart}
\usepackage{amsmath,amssymb,amsthm,graphicx}
% Retained minimal preamble for the Stacks Project chapter more-morphisms.tex (More on Morphisms).
% Verbatim excerpt of the chapter preamble: Xy-pic setup (native preamble lines 14-18) and the native
% theorem declarations with the mathematical macros (native preamble lines 154-217). Copied mechanically;
% nothing added, removed or redefined. The wrapper class and the amsmath/amssymb/amsthm/graphicx base
% packages are supplied by the assembler, so the class/externaldocument part of the native preamble is omitted.
% For commutative diagrams we use Xy-pic
\usepackage[all]{xy}

% We use 2cell for 2-commutative diagrams.
\xyoption{2cell}
% Theorem environments.
%
\theoremstyle{plain}
\newtheorem{theorem}[subsection]{Theorem}
\newtheorem{proposition}[subsection]{Proposition}
\newtheorem{lemma}[subsection]{Lemma}

\theoremstyle{definition}
\newtheorem{definition}[subsection]{Definition}
\newtheorem{example}[subsection]{Example}
\newtheorem{exercise}[subsection]{Exercise}
\newtheorem{situation}[subsection]{Situation}

\theoremstyle{remark}
\newtheorem{remark}[subsection]{Remark}
\newtheorem{remarks}[subsection]{Remarks}

\numberwithin{equation}{subsection}

% Macros
%
\def\lim{\mathop{\mathrm{lim}}\nolimits}
\def\colim{\mathop{\mathrm{colim}}\nolimits}
\def\Spec{\mathop{\mathrm{Spec}}}
\def\Hom{\mathop{\mathrm{Hom}}\nolimits}
\def\Ext{\mathop{\mathrm{Ext}}\nolimits}
\def\SheafHom{\mathop{\mathcal{H}\!\mathit{om}}\nolimits}
\def\SheafExt{\mathop{\mathcal{E}\!\mathit{xt}}\nolimits}
\def\Sch{\mathit{Sch}}
\def\Mor{\mathop{\mathrm{Mor}}\nolimits}
\def\Ob{\mathop{\mathrm{Ob}}\nolimits}
\def\Sh{\mathop{\mathit{Sh}}\nolimits}
\def\NL{\mathop{N\!L}\nolimits}
\def\CH{\mathop{\mathrm{CH}}\nolimits}
\def\proetale{{pro\text{-}\acute{e}tale}}
\def\etale{{\acute{e}tale}}
\def\QCoh{\mathit{QCoh}}
\def\Ker{\mathop{\mathrm{Ker}}}
\def\Im{\mathop{\mathrm{Im}}}
\def\Coker{\mathop{\mathrm{Coker}}}
\def\Coim{\mathop{\mathrm{Coim}}}

% Boxtimes
%
\DeclareMathSymbol{\boxtimes}{\mathbin}{AMSa}{"02}

%
% Macros for moduli stacks/spaces
%
\def\QCohstack{\mathcal{QC}\!\mathit{oh}}
\def\Cohstack{\mathcal{C}\!\mathit{oh}}
\def\Spacesstack{\mathcal{S}\!\mathit{paces}}
\def\Quotfunctor{\mathrm{Quot}}
\def\Hilbfunctor{\mathrm{Hilb}}
\def\Curvesstack{\mathcal{C}\!\mathit{urves}}
\def\Polarizedstack{\mathcal{P}\!\mathit{olarized}}
\def\Complexesstack{\mathcal{C}\!\mathit{omplexes}}
% \Pic is the operator that assigns to X its picard group, usage \Pic(X)
% \Picardstack_{X/B} denotes the Picard stack of X over B
% \Picardfunctor_{X/B} denotes the Picard functor of X over B
\def\Pic{\mathop{\mathrm{Pic}}\nolimits}
\def\Picardstack{\mathcal{P}\!\mathit{ic}}
\def\Picardfunctor{\mathrm{Pic}}
\def\Deformationcategory{\mathcal{D}\!\mathit{ef}}
\newenvironment{reference}{}{}
\newenvironment{slogan}{}{}
\newenvironment{history}{}{}
\usepackage{cleveref}
\crefname{originallemma}{original lemma}{original lemmas}
\Crefname{originallemma}{Original lemma}{Original lemmas}
\makeatletter
\newlabel{lemma-lifting-along-artinian-at-point}{{lemma-lifting-along-artinian-at-point}{source}{}{sourcepointer-1}{}}
\newlabel{lemma-lifting-along-artinian-at-point@cref}{{[originallemma][0][]lemma-lifting-along-artinian-at-point}{[1][0][]source}{}{}{}}
\makeatother
\crefname{originallemma}{original lemma}{original lemmas}
\Crefname{originallemma}{Original lemma}{Original lemmas}
\makeatletter
\newlabel{lemma-map-approximation}{{lemma-map-approximation}{source}{}{sourcepointer-2}{}}
\newlabel{lemma-map-approximation@cref}{{[originallemma][0][]lemma-map-approximation}{[1][0][]source}{}{}{}}
\makeatother
\crefname{originallemma}{original lemma}{original lemmas}
\Crefname{originallemma}{Original lemma}{Original lemmas}
\makeatletter
\newlabel{more-algebra-lemma-flat-unramified}{{0AGX}{source}{}{sourcepointer-3}{}}
\newlabel{more-algebra-lemma-flat-unramified@cref}{{[originallemma][0][]0AGX}{[1][0][]source}{}{}{}}
\makeatother
\crefname{originallemma}{original lemma}{original lemmas}
\Crefname{originallemma}{Original lemma}{Original lemmas}
\makeatletter
\newlabel{morphisms-lemma-etale-smooth-unramified}{{02GK}{source}{}{sourcepointer-4}{}}
\newlabel{morphisms-lemma-etale-smooth-unramified@cref}{{[originallemma][0][]02GK}{[1][0][]source}{}{}{}}
\makeatother
\title{An isomorphism of complete local rings is approximated by a diagram of elementary etale neighbourhoods}
\author{The Stacks Project Authors}
\date{stacks-project:more-morphisms.tex:local-snapshot-20261011; GNU FDL 1.2 or later}
\begin{document}
\maketitle
\noindent\textit{Editorial note.} Statement span is lines 11070-11095 of sources/stacks-more-morphisms/pinned-original-source.tex; proof span is lines 11097-11141, the complete author proof (43 body lines).\par\medskip
\noindent\textit{Editorial note.} Navigation stubs: 4 cited-but-not-included labels are supplied as non-load-bearing pointers with their pinned tags where the snapshot tag index has them.\par\medskip
\noindent\textit{Editorial note.} extra\_packages is empty on purpose: the retained preamble already loads xy with options and a second declaration clashes in the assembler wrapper.\par\medskip
\noindent\textit{Editorial note.} Text is retained verbatim, including the author's own typos and placeholders.\par\medskip
\noindent\textit{Allowed standard prerequisites (editorial).}\begin{itemize}
\item Explicit prerequisite boundary. G-rings, complete local rings, elementary etale neighbourhoods, and approximation of maps between completions; the construction is retained.
\end{itemize}\medskip
\ifdefined\hypertarget\hypertarget{sourcepointer-1}{}\fi
\noindent\textit{Source reference.} Original lemma lemma-lifting-along-artinian-at-point: sources/stacks-more-morphisms/pinned-original-source.tex, source lines 3575--3600. Cited by the retained statement or proof but not part of the collected spans; supplied only as a navigation pointer. The argument does not depend on its content.\par\medskip
\ifdefined\hypertarget\hypertarget{sourcepointer-2}{}\fi
\noindent\textit{Source reference.} Original lemma lemma-map-approximation: sources/stacks-more-morphisms/pinned-original-source.tex, source lines 11017--11039. Cited by the retained statement or proof but not part of the collected spans; supplied only as a navigation pointer. The argument does not depend on its content.\par\medskip
\ifdefined\hypertarget\hypertarget{sourcepointer-3}{}\fi
\noindent\textit{Source reference.} Original lemma 0AGX: sources/stacks-more-morphisms/pinned-original-source.tex, source lines 11141--11141. Cited by the retained statement or proof but not part of the collected spans; supplied only as a navigation pointer. The argument does not depend on its content.\par\medskip
\ifdefined\hypertarget\hypertarget{sourcepointer-4}{}\fi
\noindent\textit{Source reference.} Original lemma 02GK: sources/stacks-more-morphisms/pinned-original-source.tex, source lines 11141--11141. Cited by the retained statement or proof but not part of the collected spans; supplied only as a navigation pointer. The argument does not depend on its content.\par\medskip
\begin{lemma}
\label{lemma-isomorphism-approximation}
Let $S$ be a locally Noetherian scheme. Let $X$, $Y$ be
schemes locally of finite
type over $S$. Let $x \in X$ and $y \in Y$ be points lying over the
same point $s \in S$. Assume $\mathcal{O}_{S, s}$ is a G-ring.
Assume we have an $\mathcal{O}_{S, s}$-algebra isomorphism
$$
\varphi : \mathcal{O}_{Y, y}^\wedge \longrightarrow \mathcal{O}_{X, x}^\wedge
$$
between the complete local rings. Then for every $N \geq 1$
there exists morphisms
$$
(X, x) \leftarrow (U, u) \rightarrow (Y, y)
$$
of pointed schemes over $S$ such that both arrows define elementary
\'etale neighbourhoods and such that the diagram
$$
\xymatrix{
& \mathcal{O}_{U, u}^\wedge \\
\mathcal{O}_{Y, y}^\wedge \ar[rr]^\varphi \ar[ru] & &
\mathcal{O}_{X, x}^\wedge \ar[lu]
}
$$
commutes modulo $\mathfrak m_u^N$.
\end{lemma}

\begin{proof}
We may assume $N \geq 2$. Apply Lemma \ref{lemma-map-approximation} to get
$(U, u) \to (X, x)$ and $f : (U, u) \to (Y, y)$.
We claim that $f$ is \'etale at $u$ which will finish the proof.
In fact, we will show that the induced map
$\mathcal{O}_{Y, y}^\wedge \to \mathcal{O}_{U, u}^\wedge$
is an isomorphism. Having proved this,
Lemma \ref{lemma-lifting-along-artinian-at-point}
will show that $f$ is smooth at $u$ and of course
$f$ is unramified at $u$ as well, so
Morphisms, Lemma \ref{morphisms-lemma-etale-smooth-unramified}
tells us $f$ is \'etale at $u$.
For a local ring $(R, \mathfrak m)$ we set
$\text{Gr}_\mathfrak m(R) =
\bigoplus_{n \geq 0} \mathfrak m^n/\mathfrak m^{n + 1}$.
To prove the claim we look at the induced diagram
of graded rings
$$
\xymatrix{
& \text{Gr}_{\mathfrak m_u}(\mathcal{O}_{U, u}) \\
\text{Gr}_{\mathfrak m_y}(\mathcal{O}_{Y, y}) \ar[rr]^\varphi \ar[ru] & &
\text{Gr}_{\mathfrak m_x}(\mathcal{O}_{X, x}) \ar[lu]
}
$$
Since $N \geq 2$ this diagram is actually commutative as the
displayed graded algebras are generated in degree $1$!
By assumption the lower arrow is an isomorphism.
By More on Algebra, Lemma \ref{more-algebra-lemma-flat-unramified}
(for example) the map
$\mathcal{O}_{X, x}^\wedge \to \mathcal{O}_{U, u}^\wedge$
is an isomorphism and hence the north-west arrow
in the diagram is an isomorphism. We conclude that
$f$ induces an isomorphism
$\text{Gr}_{\mathfrak m_x}(\mathcal{O}_{X, x}) \to
\text{Gr}_{\mathfrak m_y}(\mathcal{O}_{U, u})$.
Using induction and the short exact sequences
$$
0 \to \text{Gr}^n_{\mathfrak m}(R) \to R/\mathfrak m^{n + 1} \to
R/\mathfrak m^n \to 0
$$
for both local rings we conclude (from the snake lemma)
that $f$ induces isomorphisms
$\mathcal{O}_{Y, y}/\mathfrak m_y^n \to \mathcal{O}_{U, u}/\mathfrak m_u^n$
for all $n$ which is what we wanted to show.
\end{proof}

\end{document}
